% Part: incompleteness
% Chapter: arithmetization-syntax
% Section: introduction

\documentclass[../../../include/open-logic-section]{subfiles}

\begin{document}

\olfileid{inc}{art}{int}
\olsection{పరిచయం}

గణనీయతను తర్కశాస్త్రంతో అనుసంధానించడానికి, తర్కశాస్త్రపు వస్తువులు (సంకేతాలు,
పదాలు, \tetoken{సూత్రాలు}{!!{formula}s}, \tetoken{వ్యుత్పత్తులు}{!!{derivation}s}),
వాటిపై క్రియలు, వాటి ధర్మాలు, సంబంధాల గురించి గణనాత్మకంగా పరిగణించదగిన విధంగా
మాట్లాడే మార్గం కావాలి. సంకేతాలు, సంకేతాల క్రమాలు, వాటి నుంచి నిర్మించిన ఇతర
వస్తువులపై గణనీయ ప్రమేయాలు, సంబంధాలను పరిగణించడం ద్వారా దీన్ని నేరుగా
చేయవచ్చు. తార్కిక వాక్యనిర్మాణపు వస్తువులన్నీ పరిమితమైనవీ,
\tetoken{లెక్కించదగిన}{!!a{enumerable}} సంకేత సమితుల నుంచి నిర్మించినవీ కనుక,
కొన్ని గణనా నమూనాలకు ఇది సాధ్యమే. కానీ పునరావృత్త ప్రమేయాల వంటి ఇతర గణనా
నమూనాలు సంఖ్యలకూ, వాటి సంబంధాలు, ప్రమేయాలకూ పరిమితమవుతాయి. అంతేకాక,
చివరికి కొన్ని సిద్ధాంతాల లోపలే, ముఖ్యంగా అంకగణిత భాషలో సూత్రీకరించిన
సిద్ధాంతాలలో, వాక్యనిర్మాణాన్ని పరిగణించగలగాలి. ఈ సందర్భాలలో వాక్యనిర్మాణాన్ని
\emph{అంకగణితీకరించడం}, అంటే వాక్యనిర్మాణ వస్తువులు, వాటిపై క్రియలు, వాటి
సంబంధాలను వరుసగా సంఖ్యలు, అంకగణిత ప్రమేయాలు, అంకగణిత సంబంధాలుగా
ప్రాతినిధ్యం చేయడం అవసరం. లైబ్‌నిట్స్ కాలానికి చెందిన ఈ ఆలోచన,
వాక్యనిర్మాణ వస్తువులకు సంఖ్యలను కేటాయించడమే.

సంకేతాలకు వాటి ``సంకేతసంఖ్యలు''గా సంఖ్యలను కేటాయించడం తులనాత్మకంగా
సూటియైన పని. కొన్ని సంకేతాలు కొద్దిగా సవాలు చేస్తాయి; ఉదాహరణకు, అనంతంగా
ఎన్నో \tetoken{చరాలు}{!!{variable}s} ఉంటాయి, ప్రతి స్థానసంఖ్య~$n$కూ అనంతంగా
ఎన్నో \tetoken{ప్రమేయ సంకేతాలు}{!!{function}s} కూడా ఉంటాయి. అయినా,
ఉదాహరణకు~$\Obj v_2$, $\Obj v_3$లకు వేర్వేరు సంకేతసంఖ్యలు వచ్చేలా సంకేతాలకు
క్రమబద్ధంగా సంఖ్యలను కేటాయించడం సాధ్యమే. సంకేతాల క్రమాలు (పదాలు,
\tetoken{సూత్రాలు}{!!{formula}s} వంటివి) ఇంకా పెద్ద సవాలు. కానీ సంఖ్యల క్రమాలను
పూర్తిగా అంకగణిత పద్ధతిలో (ఉదాహరణకు, క్రమాల ప్రధాన-సంఖ్యా ఘాత సంకేతీకరణతో)
పరిగణించగలిగితే, ఒక్కొక్క సంకేతపు సంకేతీకరణను సంకేతాల క్రమాలకు, తరువాత
\tetoken{సూత్రాల}{!!{formula}s} క్రమాలు లేదా ఇతర అమరికలకు---
\tetoken{వ్యుత్పత్తుల}{!!{derivation}s} వంటి వాటికి---విస్తరించవచ్చు. ఈ
విస్తరిత సంకేతీకరణను ``గ్యోడెల్ సంఖ్యీకరణ'' అంటారు. ప్రతి పదానికీ,
\tetoken{సూత్రానికీ}{!!{formula}}, \tetoken{వ్యుత్పత్తికీ}{!!{derivation}}
ఒక గ్యోడెల్ సంఖ్య కేటాయించబడుతుంది.

సంకేతాల క్రమాలను వాటి సంకేతసంఖ్యల క్రమాలుగా సంకేతీకరించి, గణనీయ
ప్రమేయాలతో పరిగణించగల క్రమ-సంకేతీకరణ వ్యవస్థను ఎంచుకోవడం వల్ల, గ్యోడెల్
సంఖ్యలను కూడా గణనీయ ప్రమేయాలతో పరిగణించవచ్చు. ఆచరణలో సంబంధిత ప్రమేయాలన్నీ
ఆదిమ పునరావృత్త ప్రమేయాలే అవుతాయి. ఉదాహరణకు, ఒక క్రమం పొడవును గణించడం,
క్రమపు సంకేతసంఖ్య నుంచి దాని $i$-వ మూలకాన్ని గణించడం రెండూ ఆదిమ పునరావృత్తం.
ఆ క్రమాన్ని సంకేతీకరించే సంఖ్య, ఉదాహరణకు, \tetoken{ఒక సూత్రం}{!!a{formula}}
~$!A$ యొక్క గ్యోడెల్ సంఖ్య అయితే, \tetoken{ఒక సూత్రం}{!!a{formula}} పొడవునూ,
\tetoken{ఒక సూత్రం}{!!a{formula}}లోని $i$-వ సంకేతపు (సంకేతసంఖ్యనూ)~$!A$
గ్యోడెల్ సంఖ్య నుంచి గణించవచ్చని వెంటనే తెలుస్తుంది. అయితే, ఒక సంఖ్య
సరిగా నిర్మితమైన పదపు గ్యోడెల్ సంఖ్య కావడం లేదా సరైన
\tetoken{వ్యుత్పత్తి}{!!{derivation}} కావడం ఆదిమ పునరావృత్త ధర్మమని
నిరూపించడం కొంచెం కష్టం. ఆసక్తిగల క్రమాలు (పదాలు,
\tetoken{సూత్రాలు}{!!{formula}s}, \tetoken{వ్యుత్పత్తులు}{!!{derivation}s})
ఆగమనాత్మకంగా నిర్వచించబడతాయి కనుక అది అయినా సాధ్యమే.

ఉదాహరణగా ప్రతిస్థాపన క్రియను పరిగణించండి. $!A$ ఒక సూత్రం, $x$ ఒక చరం,
$t$ ఒక పదం అయితే, $\Subst{!A}{t}{x}$ అనేది~$!A$లోని~$x$ ప్రతి స్వేచ్ఛా
ఘటనను~$t$తో మార్చడం వల్ల వచ్చే ఫలితం. ఇప్పుడు $!A$, $x$, $t$లకు వరుసగా
$k$, $l$, $m$ అనే గ్యోడెల్ సంఖ్యలు కేటాయించామని అనుకుందాం. అదే పద్ధతి
$\Subst{!A}{t}{x}$కు~$n$ అనే గ్యోడెల్ సంఖ్యను కేటాయిస్తుంది. $k$, $l$, $m$లను
$n$కు పంపే ఈ ప్రతిబింబమే ప్రతిస్థాపన క్రియకు అంకగణిత సదృశం. ప్రతిస్థాపన
క్రియ $!A$, $x$, $t$లను~$\Subst{!A}{t}{x}$కు పంపినప్పుడు, అంకగణితీకరించిన
ప్రతిస్థాపన ప్రమేయం గ్యోడెల్ సంఖ్యలు $k$, $l$, $m$లను గ్యోడెల్ సంఖ్య~$n$కు
పంపుతుంది. ఈ ప్రమేయం ఆదిమ పునరావృత్తమని చూస్తాం.
\sourcecorrection{OLTEART-001}{ఆంగ్ల మూలంలోని ``the arithmetized substitution functions maps''లో బహువచన కర్తకు ఏకవచన క్రియ ఉంది. సందర్భం నిర్దేశించే ఏకవచన అంకగణితీకరించిన ప్రతిస్థాపన ప్రమేయాన్ని, దానికి సరిపడే క్రియతో ఇచ్చాం.}

వాక్యనిర్మాణపు అంకగణితీకరణ కేవలం అమూర్త ఆసక్తికే పరిమితం కాదు. భాషల
గురించి ``మాట్లాడే'' యంత్రాంగం సహజంగా లేని అంకగణిత భాష వంటి భాషలే
వ్యక్తీకరణల సంక్లిష్ట ధర్మాలను సూత్రీకరించగలవనే విషయం మొదట సాధారణం కాని
అంతర్దృష్టి. తరువాత, పియానో అంకగణితం వంటి ఈ భాషలోని ఒక సిద్ధాంతం తన
సొంత భాష గురించి (ఉదాహరణకు, \tetoken{వాక్యాలు}{!!{sentence}s}
నిరూపించదగినవా లేదా సత్యమా అనే విషయం సహా) ఏమి \emph{నిరూపించగలదో}
అడగడానికి చిన్న అడుగు మాత్రమే మిగులుతుంది. ఇది గ్యోడెల్ (నిరూపించలేమనే
ఫలితాల గురించి), టార్స్కీ (సత్యపు అనిర్వచనీయత గురించి) ప్రసిద్ధ పరిమితి
సిద్ధాంతాలకు నడిపిస్తుంది. అయితే వాక్యనిర్మాణాన్ని అంకగణితీకరించే పద్ధతి,
సిద్ధాంతాల గణనాత్మక శక్తి లేదా వేర్వేరు గణనా నమూనాల మధ్య సంబంధం వంటి
గణనీయతా సిద్ధాంత ఫలితాలను నిరూపించడానికీ ముఖ్యమే. వాక్యనిర్మాణపు
అంకగణితీకరణ ఇతర వస్తువులు, ధర్మాలను అంకగణితీకరించడానికి ఒక నమూనాగా
పనిచేస్తుంది. ఉదాహరణకు, స్థితివిన్యాసాలు, గణనలు (ట్యూరింగ్ యంత్రాలవి
అనుకుందాం) కూడా ఇలాగే అంకగణితీకరించవచ్చు. దానివల్ల ఒక నమూనాలోని
గణనలను (ఉదాహరణకు ట్యూరింగ్ యంత్రాలు) మరొక నమూనాలో (ఉదాహరణకు పునరావృత్త
ప్రమేయాలు) అనుకరించడం సాధ్యమవుతుంది.

\end{document}
